\section*{Chapter \ref{ch:in:a-map-of-the-industry} --- A Map of the Industry}

\begin{solution}{exo:in:a-map-of-the-industry:1}
Market making for securities: 523150; speculators for own account: 523910; an investment adviser: 523940. The 2017 codes
separated a dealer (523110) from a broker (523120), and portfolio management (523920) from investment advice (523930);
523910 is unchanged.
\end{solution}

\begin{solution}{exo:in:a-map-of-the-industry:2}
$2\,891/3\,249=89.0\%$ of broker-dealers are small; of the proprietary-trading and market-making segments $93/103=90.3\%$.
\end{solution}

\begin{solution}{exo:in:a-map-of-the-industry:3}
Index fund: clients' capital, paid a fee on assets, holds for years, sells index exposure; in-house equity team of a
pension fund: beneficiaries' capital, paid from the fund's budget, holds for years, manages the fund's own assets; crypto
exchange: shareholders' capital (not at risk in trades it matches), paid transaction and listing fees, holds no position
by design, sells a venue.
\end{solution}

\begin{solution}{exo:in:a-map-of-the-industry:4}
It leaves 16 unknown: only the five issuers carry an SIC code on EDGAR, and of those it gets three right (the bank, the
exchange and the listed market maker) and two wrong (the investment bank and the asset manager under code 6211).
\end{solution}

\begin{solution}{exo:in:a-map-of-the-industry:5}
Broker-dealers: $(3\,249-3\,943)/3\,943=-17.6\%$. Adviser-only firms: $(32\,090-28\,712)/28\,712=+11.8\%$.
\end{solution}

\begin{solution}{exo:in:a-map-of-the-industry:6}
The forms rule reads every holdings-report filer as a systematic fund: it gets the group right for the two multi-manager
platforms (Millennium, Balyasny) but not their business model, so 15 right groups become 13 right models.
\end{solution}

\begin{solution}{exo:in:a-map-of-the-industry:7}
The asset manager moves to the right group: 19 of 21. The investment bank also files holdings reports (its asset
management arm) and carries SIC 6211, so the new rule calls it an asset manager, still wrong; the quantitative trading
firm whose entity files only holdings reports stays wrong.
\end{solution}

\begin{solution}{exo:in:a-map-of-the-industry:8}
The sample was chosen: its eleven broker-dealers were all market makers, so the rule ``broker-dealer means principal
trading'' was never tested on agency brokers, banks' dealer subsidiaries or retail brokers, which file the same forms.
Most of the industry's firms are small advisers or broker-dealers not in the sample. The 86\% is an in-sample score on
a convenience sample, not an accuracy.
\end{solution}

\begin{solution}{pb:in:a-map-of-the-industry:1}
\begin{enumerate}
\item Whose capital, who pays, holding period, product and markets, size; the first three define the business model.
\item Market maker: owners' capital, paid by the spread and rebates, seconds to hours. Systematic fund: investors'
capital, paid fees on assets and profits, minutes to months.
\item A proprietary trading firm that trades only its own capital, on exchanges and electronic venues, mostly as a
liquidity provider, with no clients' money; the associations add the self-description as independent market makers and
the rule that members trade their own capital.
\item Because a group can own businesses of several models, and the job belongs to one of them.
\item Principal trading 3, investment manager 5, bank 1, infrastructure 4: thirteen.
\item 523150 and 523910.
\item 523110 and 523120 into 523150; 523920 and 523930 into 523940.
\item The SEC assigns SIC codes to issuers of registered securities; a private broker-dealer is not one.
\item Goldman Sachs (bank), BlackRock (asset manager), Virtu Financial (market maker).
\item 3 correct; 16 unknown; 2 wrong.
\item Broker-dealer reports (a registered broker-dealer), 13F holdings (a manager with at least \$100 million of 13(f)
securities), issuer reports (registered securities outstanding), ATS-N (an operator of an NMS stock alternative trading
system).
\item 15 of 21 groups; 13 of 21 models.
\item Both file holdings reports and neither is a broker-dealer or an issuer.
\item 18 of 21 groups; 16 of 21 models.
\item The trading firm whose chosen entity is not its broker-dealer (read as a fund manager); the investment bank and the
asset manager (both under SIC 6211, read as dealers).
\item 3\,249 broker-dealers; 103 in the proprietary-trading and market-making segments.
\item It counts registered representatives, who deal with the public, not engineers, researchers or traders who do not.
\item SIC alone 14\% (3 of 21); forms alone 71\% (15 of 21); forms then SIC 86\% (18 of 21).
\item Step~1, listing the group's legal entities.
\item The record that decided the answer and the rule's known blind spots, so that the student can check them; and that
the tool was scored on a hand-picked sample.
\end{enumerate}
\end{solution}

\begin{solution}{iq:in:a-map-of-the-industry:1}
Whose capital is at risk (owners' against investors'), who pays the firm (the spread against fees), how long positions
are held (seconds against weeks or months). For the employee: pay comes from trading profit, not fees; feedback is fast
or slow; there are no investors to report to, or there are.
\iqlookfor{the business model, not the asset class.}
\end{solution}

\begin{solution}{iq:in:a-map-of-the-industry:2}
It decides what the trader is paid for: capturing spread with little risk (market maker), serving clients well (bank),
or earning a return on investors' money after fees (fund). It also decides what is a good day.
\iqlookfor{the link from revenue source to incentives.}
\end{solution}

\begin{solution}{iq:in:a-map-of-the-industry:3}
It is a registered broker-dealer, its audited statement of financial condition and notes are public, and it is subject
to net capital rules. It does not say whether it is a market maker, an agency broker or a bank's dealer, nor how many
people work there; its periodic FOCUS reports are confidential.
\iqlookfor{knowing what a public record proves and what it does not.}
\end{solution}

\begin{solution}{iq:in:a-map-of-the-industry:4}
Count firms: FINRA's segments list 103 proprietary-trading and market-making broker-dealers at the end of 2024; add some
50 to 100 futures-only \omterm{def:in:a-map-of-the-industry:ptf}{principal trading firms} (an assumption). Size them by band (assumptions): 80 firms of 20--100
people, 20 of 100--500, 10 of 500--3\,000. That gives $80\times20+20\times100+10\times500=8\,600$ to
$80\times100+20\times500+10\times3\,000=48\,000$: about 10\,000 to 50\,000, dominated by the ten largest firms.
\iqlookfor{a decomposition, stated assumptions, and a range driven by the largest firms.}
\end{solution}

\begin{solution}{iq:in:a-map-of-the-industry:5}
The bank's shareholders, inside a regulated balance sheet: limits come from capital rules (risk-weighted assets,
leverage) and the bank's risk appetite. The market maker's owners: limits come from its own capital, its clearing firms'
and prime brokers' margin, and its own risk limits.
\iqlookfor{capital source and the binding constraint for each.}
\end{solution}

\begin{solution}{iq:in:a-map-of-the-industry:6}
Not until it is scored on entities it was not built from: the sample is small, hand-picked and labelled by the builder,
so the score is in-sample and biased towards the cases the rules were written for. Draw a random sample of filers, label
them blind from their own statements, report per-class accuracy with intervals (18 of 21 has a wide interval), and look
at the classes the sample never contained.
\iqlookfor{selection bias, out-of-sample evaluation and uncertainty on a small count.}
\end{solution}
